% GENERATED FILE. Edit canonical lessons, then run npm run generate:books.
% canonical-source-hash: fnv1a64:3f3b37cfcaed55f5
% canonical-lessons: HI-C307-phasna, HI-C307-pukarna, HI-C307-sannata, HI-C307-nikalna, HI-C307-behal

\chapter{Trapped}
\label{ch:hi-trapped}
\hlchaptermodality{307}

\begin{chapteropening}
\textbf{By the end of this chapter:} \emph{I can say who was trapped and how they called for help.}

Complete the last lesson of chapter 307: I can say who was trapped and how they called for help.
\end{chapteropening}

\section[\texorpdfstring{\hi{फँसना} (phãsnā)}{फँसना (phãsnā)}]{\hi{फँसना} (phãsnā) — to be trapped; to be stuck}

\label{lesson:HI-C307-phasna}



\begin{quote}
Before the new one: say the Hindi for an epidemic, then the Hindi for cholera.
\end{quote}

\subsection*{You'll want to know: \hi{फँसना}}
\textbf{\hi{फँसना}} — \emph{phãsnā} — \textquotedblleft{}to be trapped; to be stuck\textquotedblright{}.

\begin{quote}
\hi{बाढ़} \hi{में} \hi{कई} \hi{लोग} \hi{छतों} \hi{पर} \hi{फँस} \hi{गए।}
\end{quote}

\begin{quote}
\emph{bāṛh meṁ kaī log chatoṁ par phãs gae.}
\end{quote}

\begin{quote}
In the flood many people were stranded on rooftops.
\end{quote}

Past (he …): \hi{वह} \hi{आदमी} \hi{फँस} \hi{गया}, \emph{vah ādmī phãs gayā}.

\begin{cousinweb}
\textbf{\hi{फँसना}} (\emph{phãsnā}) is set by CDIAL under Sanskrit \emph{spāśayati} (CDIAL 13814), of the family of \emph{pāśa}, a noose. You \textbf{\hi{फँसते} \hi{हो}} as in a net: in a flood, in traffic, or in trouble.
\end{cousinweb}

\begin{grammarlens}[title={What you've built}]
One more word for telling what happened.
\end{grammarlens}

\subsection*{Your turn}
\begin{itemize}
\raggedright
  \item \emph{Say it:} \emph{phãsnā}
  \item \emph{Say it:} \emph{phãsnā}, once more
  \item \emph{Say it:} \emph{bāṛh meṁ kaī log chatoṁ par phãs gae.} — In the flood many people were stranded on rooftops.
  \item \emph{Recall:} say the Hindi for an epidemic, then the Hindi for cholera, then say \emph{phãsnā} again
\end{itemize}

\begin{tcolorbox}[breakable,colback=teal!4,colframe=teal!35!black,title={Before you move on}]
What does \hi{फँसना} mean? (\textquotedblleft{}To be trapped; to be stuck\textquotedblright{}.) Say it once more.
\end{tcolorbox}
\section[\texorpdfstring{\hi{पुकारना} (pukārnā)}{पुकारना (pukārnā)}]{\hi{पुकारना} (pukārnā) — to call out (to)}

\label{lesson:HI-C307-pukarna}



\begin{quote}
Before the new one: say the Hindi for cholera, then the Hindi for to be trapped.
\end{quote}

\subsection*{You'll want to know: \hi{पुकारना}}
\textbf{\hi{पुकारना}} — \emph{pukārnā} — \textquotedblleft{}to call out (to)\textquotedblright{}.

\begin{quote}
\hi{छत} \hi{से} \hi{एक} \hi{बच्चा} \hi{मदद} \hi{के} \hi{लिए} \hi{पुकार} \hi{रहा} \hi{था।}
\end{quote}

\begin{quote}
\emph{chat se ek baccā madad ke liye pukār rahā thā.}
\end{quote}

\begin{quote}
From a rooftop a child was calling out for help.
\end{quote}

Past, with \hi{ने} (he or she …): \hi{उसने} \hi{मुझे} \hi{पुकारा}, \emph{usne mujhe pukārā}.

\begin{cousinweb}
\textbf{\hi{पुकारना}} (\emph{pukārnā}) is inherited from \emph{*pukkār}, a shout, a form CDIAL reconstructs (CDIAL 8246). It is a call meant to be heard far off: \textbf{\hi{मदद} \hi{के} \hi{लिए} \hi{पुकारना}}, to call for help.
\end{cousinweb}

\begin{grammarlens}[title={What you've built}]
One more word for telling what happened.
\end{grammarlens}

\subsection*{Your turn}
\begin{itemize}
\raggedright
  \item \emph{Say it:} \emph{pukārnā}
  \item \emph{Say it:} \emph{pukārnā}, once more
  \item \emph{Say it:} \emph{chat se ek baccā madad ke liye pukār rahā thā.} — From a rooftop a child was calling out for help.
  \item \emph{Recall:} say the Hindi for cholera, then the Hindi for to be trapped, then say \emph{pukārnā} again
\end{itemize}

\begin{tcolorbox}[breakable,colback=teal!4,colframe=teal!35!black,title={Before you move on}]
What does \hi{पुकारना} mean? (\textquotedblleft{}To call out (to)\textquotedblright{}.) Say it once more.
\end{tcolorbox}
\section[\texorpdfstring{\hi{सन्नाटा} (sannāṭā)}{सन्नाटा (sannāṭā)}]{\hi{सन्नाटा} (sannāṭā) — a dead silence}

\label{lesson:HI-C307-sannata}



\begin{quote}
Before the new one: say the Hindi for to be trapped, then the Hindi for to call out (to).
\end{quote}

\subsection*{You'll want to know: \hi{सन्नाटा}}
\textbf{\hi{सन्नाटा}} — \emph{sannāṭā} — \textquotedblleft{}a dead silence\textquotedblright{}.

\begin{quote}
\hi{रात} \hi{को} \hi{गाँव} \hi{में} \hi{सन्नाटा} \hi{था।}
\end{quote}

\begin{quote}
\emph{rāt ko gā̃v meṁ sannāṭā thā.}
\end{quote}

\begin{quote}
At night a dead silence lay over the village.
\end{quote}

Masculine.

\begin{cousinweb}
\textbf{\hi{सन्नाटा}} (\emph{sannāṭā}) is set by CDIAL under Sanskrit \emph{saṁnaṣṭa} (CDIAL 12925). It is not a calm quiet but a silence that frightens: \textbf{\hi{सन्नाटा} \hi{छा} \hi{गया}}, a hush fell.
\end{cousinweb}

\begin{grammarlens}[title={What you've built}]
One more word for telling what happened.
\end{grammarlens}

\subsection*{Your turn}
\begin{itemize}
\raggedright
  \item \emph{Say it:} \emph{sannāṭā}
  \item \emph{Say it:} \emph{sannāṭā}, once more
  \item \emph{Say it:} \emph{rāt ko gā̃v meṁ sannāṭā thā.} — At night a dead silence lay over the village.
  \item \emph{Recall:} say the Hindi for to be trapped, then the Hindi for to call out (to), then say \emph{sannāṭā} again
\end{itemize}

\begin{tcolorbox}[breakable,colback=teal!4,colframe=teal!35!black,title={Before you move on}]
What does \hi{सन्नाटा} mean? (\textquotedblleft{}A dead silence\textquotedblright{}.) Say it once more.
\end{tcolorbox}
\section[\texorpdfstring{\hi{निकालना} (nikālnā)}{निकालना (nikālnā)}]{\hi{निकालना} (nikālnā) — to take out; to pull out}

\label{lesson:HI-C307-nikalna}



\begin{quote}
Before the new one: say the Hindi for to call out (to), then the Hindi for a dead silence.
\end{quote}

\subsection*{You'll want to know: \hi{निकालना}}
\textbf{\hi{निकालना}} — \emph{nikālnā} — \textquotedblleft{}to take out; to pull out\textquotedblright{}.

\begin{quote}
\hi{लोगों} \hi{ने} \hi{बच्चे} \hi{को} \hi{पानी} \hi{से} \hi{निकाला।}
\end{quote}

\begin{quote}
\emph{logoṁ ne bacce ko pānī se nikālā.}
\end{quote}

\begin{quote}
People pulled the child out of the water.
\end{quote}

Past, with \hi{ने} (he or she …): \hi{उसने} \hi{बच्चे} \hi{को} \hi{निकाला}, \emph{usne bacce ko nikālā}.

\begin{cousinweb}
\textbf{\hi{निकालना}} (\emph{nikālnā}) goes back to Sanskrit \emph{niṣkālayati}, \textquotedblleft{}drives out\textquotedblright{} (CDIAL 7484). It is the transitive partner of \textbf{\hi{निकलना}}: you \textbf{\hi{निकलते} \hi{हो}} from a house, and you \textbf{\hi{निकालते} \hi{हो}} something from a bag.
\end{cousinweb}

\begin{grammarlens}[title={What you've built}]
One more word for telling what happened.
\end{grammarlens}

\subsection*{Your turn}
\begin{itemize}
\raggedright
  \item \emph{Say it:} \emph{nikālnā}
  \item \emph{Say it:} \emph{nikālnā}, once more
  \item \emph{Say it:} \emph{logoṁ ne bacce ko pānī se nikālā.} — People pulled the child out of the water.
  \item \emph{Recall:} say the Hindi for to call out (to), then the Hindi for a dead silence, then say \emph{nikālnā} again
\end{itemize}

\begin{tcolorbox}[breakable,colback=teal!4,colframe=teal!35!black,title={Before you move on}]
What does \hi{निकालना} mean? (\textquotedblleft{}To take out; to pull out\textquotedblright{}.) Say it once more.
\end{tcolorbox}
\section[\texorpdfstring{\hi{बेहाल} (behāl)}{बेहाल (behāl)}]{\hi{बेहाल} (behāl) — in a sorry state}

\label{lesson:HI-C307-behal}



\begin{quote}
Before the new one: say the Hindi for a dead silence, then the Hindi for to take out.
\end{quote}

\subsection*{You'll want to know: \hi{बेहाल}}
\textbf{\hi{बेहाल}} — \emph{behāl} — \textquotedblleft{}in a sorry state\textquotedblright{}.

\begin{quote}
\hi{दो} \hi{दिन} \hi{बाद} \hi{लोग} \hi{भूख} \hi{से} \hi{बेहाल} \hi{थे।}
\end{quote}

\begin{quote}
\emph{do din bād log bhūkh se behāl the.}
\end{quote}

\begin{quote}
Two days later people were in a sorry state from hunger.
\end{quote}

\begin{cousinweb}
\textbf{\hi{बेहाल}} (\emph{behāl}) is usually traced to Persian \emph{bē}, \textquotedblleft{}without\textquotedblright{}, with \emph{ḥāl}, \textquotedblleft{}a state, a condition\textquotedblright{}, an Arabic word Persian took over: without a good state, at the end of one's strength.
\end{cousinweb}

\begin{grammarlens}[title={What you've built}]
One more word for telling what happened.
\end{grammarlens}

\subsection*{Your turn}
\begin{itemize}
\raggedright
  \item \emph{Say it:} \emph{behāl}
  \item \emph{Say it:} \emph{behāl}, once more
  \item \emph{Say it:} \emph{do din bād log bhūkh se behāl the.} — Two days later people were in a sorry state from hunger.
  \item \emph{Recall:} say the Hindi for a dead silence, then the Hindi for to take out, then say \emph{behāl} again
\end{itemize}

\begin{tcolorbox}[breakable,colback=teal!4,colframe=teal!35!black,title={Before you move on}]
What does \hi{बेहाल} mean? (\textquotedblleft{}In a sorry state\textquotedblright{}.) Say it once more.
\end{tcolorbox}
